% Section 7: Conclusion
\section{Conclusion}
\label{sec:conclusion}

We began by asking whether a model preferred by humans is also preferred after length debiasing
--- and found, across 223 models, an unambiguous yes with a partial correlation of 0.985 that
exceeds the already-high bivariate agreement. Two evaluation systems agree on model capability
ordering at a level the LLM evaluation community has not previously quantified.

Our contributions are: (1) $r_{\mathrm{partial}} = 0.985$ ($p = 1.69 \times 10^{-170}$,
$N=223$) establishing population-scale capability-LC alignment; (2) capability-verbosity
dominance $4.88{:}1$ in standardized OLS, with negative verbosity coefficient confirming LC
correction functionality; (3) FWL verification within 1.1pp ruling out artifact; (4) monotonic
quartile ordering with $\varepsilon^2 = 0.883$; and (5) a $10\times$ DV-choice effect-size
lesson favoring \texttt{LC\_winrate} over $\Delta$.

Future work: cross-framework replication (MT-Bench, Chatbot Arena, $N \geq 100$ overlap);
family-stratified analysis; longitudinal extension to post-2024 frontier models; and
Pearson-based FWL exact equality verification.

The LC correction is not merely a statistical adjustment --- it is a validated
capability-preserving transformation. Future evaluation research can build on this foundation
with confidence.
